\chapter{Unified Barrick Valuation Experiment and Validation}
\label{ch:unified-experiment}

This chapter reports one controlled comparison of four gold-price engines
under a common operating and DCF contract, then evaluates its engineering,
numerical, accounting and external validity. Chapter~9 remains the home of
option calibration, residual geometry and out-of-sample repricing; those
results are referenced rather than repeated here.

\section{End-to-end experiment contract}

The numerical experiment values the eight-mine gold forecast of Chapter~10
from the valuation time onwards. Section~\ref{sec:copper-extension} reports a
separate extension with three attributable copper operations and a paired
comparison of explicit flows.

Chapter~11 defined the accounting, discounting and closure rules. The
integrated experiment is a sequence of typed transformations rather than one
black box:
\begin{align}
\mathcal P_{opt}^{\mathbb Q}
&\longrightarrow
\{S_{t,j}^{GLD,(m)},v_{t,j}^{(m)},N_{t,j}^{J,(m)},
\lambda_{t,j}^{H,(m)}\},\\
&\xrightarrow{\;\mathcal A_{scenario}\;}
\{P_{t,j}^{Au,(m)}\},\\
&\xrightarrow{\;\mathcal A_{operations}\;}
\{Rev_{q,j}^{(m)},\text{component margin}_{q,j}^{(m)}\},\\
&\xrightarrow{\;\mathcal A_{accounting}\;}
\{NOPAT_{q,j}^{(m)},FCFF_{q,j}^{proxy,(m)}\},\\
&\xrightarrow{\;\mathcal A_{DCF}\;}
\{V_{op,j}^{(m)}\},
\qquad j\in\{BS,H,BP,BH\}.
\end{align}
Each arrow has its own date, unit, measure and validation rule. Team~8 governs
the option-implied gold-price engines of Chapter~9; Team~4 supplies the
operating forecasts; the unified DCF governs cash-flow construction and
discounting; Team~3 supplies simulation diagnostics. The last step stops at
an operating proxy: no equity bridge is applied.

\begin{table}[H]
\centering
\caption{Interfaces and stop conditions in the unified experiment.}
\label{tab:unified-interface-contract}
\scriptsize
\setlength{\tabcolsep}{3pt}
\renewcommand{\arraystretch}{1.10}
\begin{tabularx}{\textwidth}{@{}p{2.15cm}p{2.15cm}X X@{}}
\toprule
Stage & Source authority & Required contract & Stop condition \\
\midrule
Option scenarios & Teams 8 and 6 & Dated snapshot, rate dates, carry, parameter
objects, \(\mathbb Q\), time grid and seed. & A calibration or predictive
claim is inconsistent with Chapter~9. \\
Gold bridge & Team 8; unified adapter & Option-implied distributional shape
applied to gold USD/oz from the spot at the valuation time. & A physical
forecast or \(\mathbb Q\)-to-\(\mathbb P\) transformation is implied rather
than declared. \\
Operations & Team 4 & Production and cost units, ownership, frequency,
dependence and accounting definition. & CoS/AISC/D\&A/capex treatment or input
perimeter is ambiguous. \\
FCFF and closure & Team 5 & Taxes, reinvestment, WACC, horizon and terminal
policy. & Accounting identities, \(w>g\) or finite-resource closure fail. \\
Equity and per share & Team 5 plus corporate inputs & Net debt, minorities,
other claims, non-operating assets, diluted shares and matching listing. &
Not applied: bridge items are not reconciled to one dated filing. \\
\bottomrule
\end{tabularx}
\end{table}

\section{Controlled design: one changing layer}

Only the gold-price scenario generator changes across the four branches.
Production, unit costs, taxes, WACC shocks and terminal policy are common.
Differences among the four output distributions can therefore be attributed
to the gold-price engine within this contract. They cannot be attributed to
mine plans, extraction uncertainty or cost inflation, because those layers do
not vary. The editable values and their terminal-value effect are disclosed
once in Section~\ref{sec:executable-valuation-assumptions}.

The earlier Team~4 NSS, IV, GBM and stochastic EBITDA chain was a standalone
demonstrator and is not imported. Team~4 contributes only the operating
decomposition and forecasts; Team~8 alone generates the price paths. The
Treasury NSS term structure of 2 September 2026 is a separate, newly fitted
input.

\section{Valuation clock}

The valuation time is the 2 September 2026 close (20:00~UTC), and only flows
after it are valued. Barrick's Q1--Q2 2026 actuals are excluded: they have
been realised and belong to the balance sheet, not to future cash flows. The
quarter containing the valuation date contributes its remaining \StubDays{} of
92 days, under the registered assumption of uniform production and unit cost
within the quarter. Eighteen quarters follow, to 31 December 2030
(\HorizonYears{} years). Revenue uses the average simulated daily gold price
over each quarter's remaining interval; each flow is discounted from the
midpoint of that interval; the first-year WACC is $w_0$. The terminal value
is set at 31 December 2030 on the last four quarters.

\section{Measure, dates and simulation design}

Option calibration identifies a risk-neutral law. The scenario drift is the
integral-preserving forward curve of the Treasury fit. Read as a physical
expectation, this assumes no gold risk premium, a common assumption for an
asset with near-zero market beta that is not estimated here. The
option-implied volatility shape is carried over as a stated conditional
assumption, not a validated physical law. Expected flows are discounted at
WACC, as in an ordinary DCF; the explicit flows are also discounted at the
risk-free curve for comparison (Table~\ref{tab:unified-terminal}).

A genuine measure transformation may change variance drift, jump intensity
and mark distribution as well as price drift; Section~\ref{sec:physical-measure-bridge}
documents the cross-measure link and the separate physical forecasting
evidence. For copper, the cross-delivery futures curve imposes a scenario mean
and must not be identified with the time drift of a fixed-delivery future.

All dates are synchronous: the gold spot anchor (USD~4,391.32/oz, XAUUSD at
20:00~UTC), the GLD surface, the Treasury curve and the valuation time share 2
September 2026. GLD carries the 0.40\% sponsor fee as a continuous yield.

The engines are simulated on a daily grid with 8,192 paths and one base seed
(20260902) for all four; WACC uses its own seed and is generated once. Variance
follows a full-truncation Euler scheme with an unclipped state; the Hawkes
intensity decays exactly between days and the jump compensator uses the same
exact integral. Monte Carlo call prices at one and three years lie within
about two batch standard errors of the COS prices for every model.

\section{Results}

For \(M\) paths from model \(j\), the empirical distribution and quantiles of
the operating proxy are
\begin{equation}
\widehat F_j(x)=\frac{1}{M}\sum_{m=1}^{M}
\mathbf 1\!\left\{V_{op,j}^{(m)}\leq x\right\},
\qquad
\widehat Q_{u,j}=\inf\{x:\widehat F_j(x)\geq u\}.
\end{equation}

\begin{table}[H]
\centering
\caption{Operating proxy under the four gold-price laws, USD billion (signed terminal value).}
\label{tab:unified-quantiles}
\small
\begin{tabular}{@{}lrrrrr@{}}
\toprule
Gold engine & P10 & Median & P90 & Mean & Negative \\
\midrule
Black--Scholes/GBM & \ValBSPTen & \ValBSMedian & \ValBSPNinety & \ValBSMean & \ValBSNegative\% \\
Heston & \ValHPTen & \ValHMedian & \ValHPNinety & \ValHMean & \ValHNegative\% \\
Bates--Poisson & \ValBPTen & \ValBMedian & \ValBPNinety & \ValBMean & \ValBNegative\% \\
Full Bates--Hawkes & \ValWPTen & \ValWMedian & \ValWPNinety & \ValWMean & \ValWNegative\% \\
\bottomrule
\end{tabular}
\end{table}

\begin{figure}[H]
  \centering
  \includegraphics[width=0.85\textwidth]{\FigFigValPrimary}
  \caption{Full Bates--Hawkes operating proxy on the dated clock, 8,192 paths.
  Display limits are the 0.5th--99.5th percentiles. This is not a confidence
  interval or a physical probability law.}
  \label{fig:unified-primary}
\end{figure}

Medians range from USD~\ValMedianLow{} to \ValMedianHigh~billion across the
four engines. Negative values are not forecasts of a negative limited-liability
share price: they expose the unconstrained operating mapping.

\begin{table}[H]
\centering
\caption{Median proxy, USD billion, under the terminal conventions and with risk-free discounting of the explicit flows.}
\label{tab:unified-terminal}
\small
\begin{tabular}{@{}lrrr@{}}
\toprule
Gold engine & Signed perpetuity & Explicit, WACC & Explicit, risk-free \\
\midrule
Black--Scholes/GBM & \ValBSMedian & \ValBSExplicit & \ValBSExplicitRiskFree \\
Heston & \ValHMedian & \ValHExplicit & \ValHExplicitRiskFree \\
Bates--Poisson & \ValBMedian & \ValBExplicit & \ValBExplicitRiskFree \\
Full Bates--Hawkes & \ValWMedian & \ValWExplicit & \ValWExplicitRiskFree \\
\bottomrule
\end{tabular}
\end{table}

The terminal value carries about \TerminalShare\% of the mean; the explicit
flows to 2030 are worth USD~\ValExplicitLow--\ValExplicitHigh~billion.
Discounting the explicit flows at the risk-free curve instead of WACC raises
them by about one billion.

\section{Numerical robustness}

Doubling the paths to 16,384 or changing the seed moves single-model medians
across USD~\NumericalMedianLow--\NumericalMedianHigh~billion; for Bates--Hawkes
alone the range is \NumericalHawkesLow--\NumericalHawkesHigh. This numerical
range is as wide as the spread between the four models, so the medians do not
rank the models. Monte Carlo error and economic dispersion are distinct: the
first concerns whether an estimator has converged, the second the spread
induced by the assumed stochastic and corporate inputs.

\section{What the corrections changed}

Table~\ref{tab:unified-attribution} adds the corrections of 24 October 2026 one
at a time to the medians published on 7 September. The dated clock removes the
first-half 2026 margin that was previously discounted as a future flow and
ends the explicit horizon at 2030; between the two editions the explicit flows
fall from about USD~13 to 11~billion. The spot anchor and the GLD carry each
move the medians by less than the numerical range.

\begin{table}[H]
\centering
\caption{Median proxy, USD billion, as each correction is added.}
\label{tab:unified-attribution}
\small
\begin{tabular}{@{}lrrrr@{}}
\toprule
Step & BS/GBM & Heston & Bates & Bates--Hawkes \\
\midrule
7 September Paper & \AttLegacyBS & \AttLegacyH & \AttLegacyB & \AttLegacyW \\
+ dated clock, path average, mid-period, corrected engines & \AttClockBS & \AttClockH & \AttClockB & \AttClockW \\
+ spot anchor at the valuation time & \AttAnchorBS & \AttAnchorH & \AttAnchorB & \AttAnchorW \\
+ GLD carry $q=0.40\%$ & \AttCarryBS & \AttCarryH & \AttCarryB & \AttCarryW \\
\bottomrule
\end{tabular}
\end{table}

\begin{figure}[H]
  \centering
  \includegraphics[width=0.80\textwidth]{\FigFigCorrectionAttribution}
  \caption{Median proxy of each engine as the corrections are added.}
  \label{fig:unified-attribution}
\end{figure}

\section{Layered validation}

Validation is not one score. Source parity, implementation tests, data
integrity, calibration, prediction, simulation stability and corporate
reconciliation answer different questions.

\begin{table}[H]
\centering
\caption{Evidence hierarchy for the integrated valuation.}
\label{tab:validation-evidence-hierarchy}
\scriptsize
\setlength{\tabcolsep}{3pt}
\renewcommand{\arraystretch}{1.10}
\begin{tabularx}{\textwidth}{@{}p{2.20cm}X X X@{}}
\toprule
Layer & Accepted evidence & Still open & Allowed conclusion \\
\midrule
Source integrity & Hashes and parity outputs for the eight team modules. &
Historical runtimes, licences and excluded assets recorded in their source
audits. & The files used are stable inputs, not independently re-estimated
facts. \\
Implementation & Offline test suite; path engines checked against COS prices;
nested calibration candidates retained. & Full independent formula review and
all historical-market replay. & The tested contracts are implemented within
their tolerance. \\
Experiment identity & Same operating contract, DCF inputs and WACC paths in all
branches. & Stochastic production and cost propagation; filing reconciliation.
& Differences among branches isolate price-engine sensitivity. \\
Accounting & Discounting identities are executable. & Debt, cash, minorities,
other claims, share count and finite-mine closure. & Outputs are operating
proxies, not fair values. \\
External validity & Dates, units, carry and measure semantics are recorded. &
Physical gold-price law, realised selling prices and risk premia. & The
experiment is reproducible but not a physical gold forecast. \\
\bottomrule
\end{tabularx}
\end{table}

\section{Accounting and economic validation}

The operating vectors stay deterministic, which supports attribution but
suppresses production and cost uncertainty. The unified DCF is a common
research contract, not a filing-reconciled model: debt, cash, minority
interests, non-operating assets and other claims are not reconciled, so no
equity value or value per share is reported. The forecast also covers only the
eight-mine subset rather than all of Barrick's gold production.

Economic validation must additionally reconcile revenue to volumes and
realised prices; operating margin to EBIT; NOPAT, reinvestment and FCFF by path
and quarter; and the adjustments from enterprise to equity value described in
Chapter~11. No option-pricing fit can substitute for these controls.

\section{External validity and information limits}

GLD option values, the GLD share price, physical gold in USD/oz, Barrick's
realised selling price and mine-level cash flow are distinct objects. The run
applies an option-implied distributional shape to gold paths that start at the
spot price, without a validated physical-measure map. The information set also
excludes timestamped macroeconomic and geopolitical news; such events may
enter indirectly through prices, implied volatility and jumps, but their
causal contribution cannot be identified here.

\section{What the integrated experiment establishes}

Under one set of operating and corporate assumptions, the four gold-price
engines give medians within the Monte Carlo range of any one of them, while the
perpetuity after 2030 carries most of the value. The choice among these
option-implied laws therefore matters less for this proxy than the numerical
precision of the simulation and, above all, the terminal assumptions. What
remains open is the economic validation of the commodity, measure, operating,
accounting and terminal-closure interfaces, summarised in the conclusion that
follows.
